\documentclass[11pt]{article}
\usepackage[margin=0.88in]{geometry}
\usepackage{amsmath,amssymb,amsthm,mathtools,mathrsfs}
\usepackage{microtype}
\usepackage{enumitem}
\usepackage{booktabs}
\usepackage[hidelinks]{hyperref}
\setlist{nosep,leftmargin=*}
\newtheorem{theorem}{Theorem}[section]
\newtheorem{proposition}[theorem]{Proposition}
\newtheorem{lemma}[theorem]{Lemma}
\newtheorem{corollary}[theorem]{Corollary}
\theoremstyle{definition}
\newtheorem{definition}[theorem]{Definition}
\newtheorem{warning}[theorem]{Warning}
\theoremstyle{remark}
\newtheorem{remark}[theorem]{Remark}
\newcommand{\Tr}{\operatorname{Tr}}
\newcommand{\Ends}{\operatorname{Ends}}
\newcommand{\Sub}{\operatorname{Sub}}
\newcommand{\ii}{\partial_{\!i}}
\title{Proposed Manuscript Insertions VIII\\
Marked Pulled/Immersed Ideal Interfaces and Actual Trace Welding}
\author{Research draft for the marked-width MLC program}
\date{July 25, 2026}

\begin{document}
\maketitle

\begin{abstract}
We separate the coarse ideal marking on which a covering/immersion pair acts
from a finer response/access trace overlay.  For a finite pseudo-piece boundary
atlas we give a necessary-and-sufficient compiler for branch-adapted ideal
markings, including exact overlap and covering-refinement obstructions.  We
construct collared ideal-arc trace spaces, proper and end-proper end-germ trace
objects, and the actual one-to-many welding obtained by immersion restriction,
covering subdivision, and covering transport.  The complete marked interface
is functorial under words, normalized overlays, and certified final-root
commuting squares.  This closes a fixed-horizon direct-decoration interface
module; universal horizon realization, H2, MBTT, UPTCR, and MLC remain open.
\end{abstract}

\section{The marked correspondence}
Let
\[
 F=(f,\iota):V'\rightrightarrows V
\]
be a holomorphic covering and a holomorphic immersion.  For an ideal marking
$\mathcal I$ of $V$, the KKL action condition is
\[
 \boxed{f^*\mathcal I\text{ refines }\iota^*\mathcal I.}
 \tag{1.1}
\]
The immersion pullbacks may fail to exist or overlap, and properness or
end-properness is additional data.

\begin{definition}[Two-layer interface]
A finite marked interface is
\[
 \mathfrak I_F=(\mathcal I,\mathscr O,f^*\mathcal I,
 \iota^*\mathcal I,\mathfrak p,\mathfrak e,
 \mathcal T,\mathcal W,\mathfrak F),
\]
where $\mathcal I$ is the action marking and $\mathscr O$ is a possibly finer
trace/access/obstacle overlay.  An endpoint of $\mathscr O$ lying inside an
interval of $\mathcal I$ is response-only until it is successfully promoted by
the action-marking compiler.
\end{definition}

\begin{warning}
Trace visibility does not imply ideal-action visibility.  No response
subdivision is used as a KKL endpoint without proving (1.1).
\end{warning}

\section{Finite pulled/immersed marking compiler}
Cut one ideal-boundary component at a retained phase endpoint.  Let the target
be tiled by micro-atoms $M_0,\ldots,M_{N-1}$ and the source by
$A_0,\ldots,A_{m-1}$.  Each source atom has a covering label $c(j)$; under the
immersion it has either a proper target label $r(j)$ or the nonproper label
$\bot$.

For a connected target atom union $B$, delete every proper source atom whose
label lies outside $B$, and retain the remaining connected components that
contain a proper atom labelled in $B$.  Denote their union by $H_\iota(B)$.

\begin{theorem}[Finite immersion pullback formula]
For a boundary-monotone finite pseudo-piece atlas,
\[
 \boxed{\iota^*B=H_\iota(B).}
 \tag{2.1}
\]
\end{theorem}

\begin{proof}
A proper path subordinate to a retained component can land only on a proper
atom mapped into $B$; intervening nonproper atoms have no target ideal landing
interval.  Hence its immersed image is subordinate to $B$.  Conversely, every
excluded atom is proper outside $B$, or is separated from all $B$-labelled
atoms by such a blocker.  A proper path landing at that blocker witnesses
failure of subordination to $B$.  Maximality in the definition of
$\iota^*B$ proves the claim.
\end{proof}

A maximal nonproper source run with adjacent proper target labels $p,q$ belongs
to both pullback hulls unless $p,q$ lie in one target interval.

\begin{theorem}[Finest immersion-admissible marking]
Join the proper labels adjacent to every maximal nonproper run and take the
interval-convex transitive closure in the retained phase order.  The resulting
contiguous blocks form the unique finest target partition whose immersion
pullbacks are pairwise disjoint.  Every immersion-admissible marking is a
coarsening of this partition.
\end{theorem}

\begin{proof}
The overlap across every nonproper run forces the adjacent labels into one
interval, and connectedness forces all intervening target micro-atoms into the
same interval.  After taking the closure, a nonproper run has proper neighbors
in at most one target block.  Proper atoms of other blocks are blockers in
(2.1), so distinct pullbacks are disjoint.
\end{proof}

Let $K$ be a connected component of an immersion pullback and let
$v_-,v_+$ be its source boundary vertices.  Write $c_\partial$ for the covering
map on source vertices.

\begin{theorem}[Covering/immersion action criterion]
For an immersion-admissible target marking $\mathcal I$, condition (1.1) holds
if and only if
\[
 \boxed{c_\partial(v_-),c_\partial(v_+)\in\partial\mathcal I}
 \tag{2.2}
\]
for every connected component of every immersion pullback.
\end{theorem}

\begin{proof}
The covering marking is obtained by pulling back target endpoints.  It refines
the induced immersion marking exactly when every boundary of every immersion
interval is one of these covering-preimage endpoints.
\end{proof}

\begin{corollary}[Finite decision]
For a finite mandatory endpoint set, enumerate the coarsenings of the finest
immersion-admissible partition, retain the mandatory cuts, and test (2.2).
This returns a finest valid action marking with all proper/end-proper flags, or
an exact obstruction: a mandatory cut through an immersion-forced block, or a
covering boundary image incompatible with every permitted coarsening.
\end{corollary}

\section{Collared ideal-arc trace spaces}
Every pseudo-piece ideal tile arc lies on a universal-cover piece used in the
source construction.  For a proper finite ideal subarc $J$, choose a
piecewise-analytic collar $C_Y(J)$ bounded by $J$, two crosscuts, and one
interior cross-section.

\begin{definition}[Collared trace]
Set
\[
 \mathcal T_Y(J)=\Tr_J H^1(C_Y(J))
\]
with the minimum-energy quotient norm
\[
 \|\tau\|^2=\inf_{\Tr_Ju=\tau}
 \left(\int_{C_Y(J)}|\nabla u|^2+
 \left|\int_{K_J}u\right|^2\right),
 \tag{3.1}
\]
where $K_J$ is a fixed compact normalization set.
\end{definition}

\begin{theorem}[Ideal-arc trace theorem]
$\mathcal T_Y(J)$ is a Hilbert trace space.  In a conformal rectangle
coordinate it is the usual $H^{1/2}$ trace space, and transporting the collar
by a conformal isomorphism transports the Dirichlet seminorm isometrically.
\end{theorem}

\begin{proof}
Apply the Sobolev trace theorem to the piecewise-analytic collar and use the
minimum-energy quotient.  A Riemann map identifies the collar with a rectangle;
Dirichlet energy is conformally invariant in dimension two.  The normalization
term removes the constant kernel.
\end{proof}

For finitely many arcs, the compatible tuple space is the image of one
$H^1$ function on the union of their collars and a finite connecting core.  It
is not replaced by an arbitrary direct sum.

If a component $K$ of $\iota^*I$ is proper, retain its full compatible trace.
If it is only end-proper, retain the compatible traces on its proper end
intervals:
\[
 \mathcal T_{V'}^{\rm end}(K)
 =\mathcal T_{V'}(E_-(K),E_+(K)).
 \tag{3.2}
\]
The nonproper middle remains in the interval bridge word and is not assigned a
fictitious boundary trace.

\section{Actual one-to-many welding}
Let $I$ be a target action interval.  For each proper end arc
$E\subset K\subset\iota^*I$, immersion restriction gives
\[
 R_{\iota,E}:\mathcal T_V(I)\longrightarrow\mathcal T_{V'}(E).
\]
Because $f^*\mathcal I$ refines $\iota^*\mathcal I$, subdivide $E$ into
covering intervals $E_\ell$.  Degree-one covering transport gives
\[
 P_{f,E_\ell}:\mathcal T_{V'}(E_\ell)
 \longrightarrow\mathcal T_V(f(E_\ell)).
\]

\begin{definition}[Marked welding]
\[
 \begin{aligned}
 \mathcal W_{F,I}:={}&
 \bigoplus_{K,E,E_\ell}
 \bigl(P_{f,E_\ell}\circ\Sub_{E_\ell}\circ R_{\iota,E}\bigr).
 \end{aligned}
 \tag{4.1}
\]
The output index retains branch, component, interval, side, phase, endpoint
flags, and the proper-or-nonproper bridge word.
\end{definition}

\begin{theorem}[Actual ideal welding]
The operator (4.1) is continuous on the compatible trace spaces.  It is the
complete trace welding when $F$ is proper on $I$, and the complete end-germ
welding when $F$ is end-proper on $I$.  If $\nu(I)$ is the number of proper
lifted end arcs, then
\[
 \|\mathcal W_{F,I}\|\le\sqrt{\nu(I)}
 \tag{4.2}
\]
in direct-sum Dirichlet norm.
\end{theorem}

\begin{proof}
Each factor is an actual trace restriction or conformal transport on a proper
collar.  The action condition prevents a covering endpoint from crossing an
immersion interval.  One parent trace is copied to at most $\nu(I)$ disjoint
proper collars; conformal energy invariance gives (4.2).
\end{proof}

\begin{theorem}[Word functoriality]
For composable marked correspondences with matching intermediate action
markings,
\[
 \boxed{
 \mathcal W_{F_2\circ F_1,I}
 =\mathcal W_{F_2,F_{1*}I}\circ\mathcal W_{F_1,I},
 }
 \tag{4.3}
\]
after canonical reassociation of the labelled compatible tuple spaces.  The
pulled intervals, properness flags, bridge words, endpoint flags, and cyclic
orders compose in the same chronological order.
\end{theorem}

\begin{proof}
Both sides restrict the same finite-energy function along the same actual
immersions, subdivide by the same covering endpoints, and transport through the
same covering maps.  Function restriction and map composition are associative.
\end{proof}

\section{Repair and overlay rules}
A trace-overlay endpoint not belonging to the action marking remains
response-only.  If a later KKL theorem requires it as an ideal endpoint, rerun
the finite compiler.

\begin{definition}[Slit-repair certificate]
A finite source/target analytic crosscut pair is a slit-repair certificate if it
avoids all critical, protected, obstacle, port, and localization data; commutes
with both maps; and separates the conflicting nonproper bridge after cutting.
\end{definition}

\begin{theorem}[Certified slit repair]
Cutting along a slit-repair certificate lifts the covering and immersion to the
cut surfaces, converts the certified split into a genuine ideal endpoint, and
reduces the repaired interface to the finite compiler above.
\end{theorem}

\begin{proof}
A finite analytic cut replaces each crosscut by two proper boundary sides.  The
commuting pullback condition lifts both maps to the slit surfaces.  The former
nonproper bridge is separated at the certified side, so the new endpoint is no
longer inside an immersion-forced block.
\end{proof}

The theorem realizes a repair from a certificate; it does not construct such a
certificate for every mixed block.

\section{Normalized and final-root naturality}
The binary pseudo-piece normalization transports the root collars, action
markings, and overlay to every finite word.  Source functoriality and (4.3) give
exact normalized-overlay naturality; no physical-coordinate welding conjugacy
is omitted.

Consider an old and enlarged pair with embeddings
\[
 A:V\to V^+,
 \qquad A':V'\to V^{+'}
\]
satisfying
\[
 \boxed{A\circ f=f^+\circ A',\qquad
 A\circ\iota=\iota^+\circ A'.}
 \tag{6.1}
\]
Assume exact inverse image of all proof-used endpoint sets, transported collars,
phase and side data, and injective cyclic-order-preserving ideal boundary maps.

\begin{theorem}[Certified final-root naturality]
Under (6.1), covering and immersion pullbacks commute with enlargement, the
action condition is preserved, proper/end-proper flags are unchanged, and
\[
 \boxed{
 A_{\rm out}^{\sharp}\mathcal W_{F^+,A(I)}
 =\mathcal W_{F,I}A_{\rm in}^{\sharp}
 }
 \tag{6.2}
\]
after matching the labelled output tuple.
\end{theorem}

\begin{proof}
The first conclusions follow by taking inverse images in the two commuting
squares.  Trace restriction by $A$ followed by old immersion pullback is the
same function as enlarged immersion pullback followed by $A'$.  The covering
square gives the output identity.  Subdivision and direct sums preserve these
equalities, proving (6.2).
\end{proof}

A finite overlay refinement is natural under the same rule.  Trace-only
subdivision always commutes.  Promotion to an action endpoint is allowed only
when the compiler accepts it or a slit repair is supplied.

\section{Fixed-horizon completion}
\begin{theorem}[Fixed-horizon marked binary interface]
Fix one actual direct-decoration binary block and a finite source-authorized
branch horizon.  Freeze all mandatory ideal endpoints, trace/access/obstacle
endpoints, proper/end-proper calls, and final-root generators.  Assume every
mandatory action endpoint passes the compiler or has a slit-repair certificate,
and every enlargement generator supplies (6.1) with the stated boundary data.
Then the complete horizon has:
\begin{enumerate}
\item actual covering-pulled and immersion-pulled ideal markings;
\item proof of the action condition on every branch;
\item complete proper/end-proper and endpoint-flag records;
\item valid collared ideal-arc and end-germ trace spaces;
\item actual one-to-many welding and bridge words;
\item exact word, normalization, overlay, and final-root naturality;
\item a finite first-failure certificate otherwise.
\end{enumerate}
\end{theorem}

\begin{proof}
Run the finite marking compiler at every branch, apply certified repairs, build
the collared trace spaces, and use the welding and naturality theorems above.
The horizon is finite, so the union of all records is finite.
\end{proof}

\section{Audit and promotion boundary}
The companion compiler was checked against an independent exhaustive search of
every target endpoint subset in 1,500 randomized systems with at most seven
target micro-atoms.  It compiled 1,221 systems, obstructed 279, and had zero
false positives, false negatives, or maximal-marking mismatches.  A second
exact-integer audit checked 1,000 welding composition and enlargement squares
with zero failures.

The theorem closes a fixed-horizon marked-interface module.  It does not build
the complete mixed Wave horizon, universally realize every mandatory endpoint
or slit, bound lift multiplicity, construct a universal response/WAD/Value
quotient, prove MBTT, or prove MLC.

\section{Next obligation}
The next source theorem is
\[
 \boxed{\texttt{UNIVERSAL\_MANDATORY\_ENDPOINT\_VISIBILITY\_AND\_REPAIR}.}
\]
It must prove that every proof-required ideal endpoint in every mixed block is
compiler-visible or has a collision-free slit repair, remains natural under
final-root enlargement, has universally bounded proper-lift multiplicity, and
factors with all response/WAD/Value outputs through one uniformly finite future
quotient.

\begin{thebibliography}{9}
\bibitem{KL08}
J. Kahn and M. Lyubich,
\emph{A priori bounds for some infinitely renormalizable quadratics. II.
Decorations}, arXiv:math/0609046.
\bibitem{KKL26}
J. Kahn, A. Kapiamba, and M. Lyubich,
\emph{MLC for parabolically bounded primitive renormalization},
arXiv:2606.27272.
\end{thebibliography}
\end{document}
